% ED26-BM-RAT2021-10-3: exact Meggle 2021 numbered chapter, pp. 601--610.
\chapter{Communicative Understanding, Acceptance, and Rationality}
\label{sec:p18-rationality-10-3}

A listener can understand an assertion without accepting it, accept it
without its being true, and participate correctly in a procedure without
thereby obtaining evidence for a factual conclusion. Georg Meggle's 2021
chapter distinguishes these questions through an individualistic account
of communicative action, a collectivistic reading of speech-act theory,
and Habermas's theory of collective communicative action
\cite[pp.~601--610]{Meggle2021RationalityCommunication}.
Its subject is intentional communicative \emph{action}, not every process
that transfers information. The account below retains the chapter's
definitions, qualifications and disagreements; it does not supply one
universal mathematical or empirical theory of communication.

\section{What the individualistic account starts from}

\begin{hbdefinition}{def:p18-r103-action-notation}{Actions and attitudes in the selected account}
At a fixed time \(t\), the source writes \(D(X,f)\) for agent
\(X\)'s performing action token \(f\), \(B(X,A)\) for the agent's
strong belief or conviction that \(A\), and \(W(X,A)\) for the
agent's wanting \(A\). Primes mark a later time \(t'>t\), rather
than a derivative or a different belief strength. Its first definition is
\[
 I(X,f,A)\defeq D(X,f)\land W(X,A)
                    \land B(X,D(X,f)\leftrightarrow A).
\]
Thus an intention attributed to an actual performance combines an action,
a wanted outcome, and a believed means--end connection. This is the
source's qualitative explication, not a definition of an intention that
can exist before or without the action.

The displayed biconditional is material equivalence. Meggle says that
it should be strengthened to a subjunctive biconditional but sets that
complication aside. Note 2 prefers equivalence to mere sufficiency in
order to address why this particular action is selected. This preference
is not a derivation of a unique choice from all available alternatives:
those alternatives, costs and causal dependencies have not been encoded.
\cite[pp.~601--602; p.~609, n.~2]{Meggle2021RationalityCommunication}
\index{communication!action rationality}\index{intention!communicative}
\end{hbdefinition}

In this restricted representation success requires that the agent's
means--end belief be correct: \(D(X,f)\leftrightarrow A\) actually
holds. Since \(I(X,f,A)\) already includes the action's occurrence,
the displayed success test amounts to attaining \(A\). It does not
by itself establish that the action caused \(A\); an independently
occurring outcome can satisfy a material biconditional. False beliefs
can nevertheless make the performance intelligible in a subjective
means--end explanation. A primary aim is an aim for which the other
aims are means, if such an aim exists; the chapter does not assume
that every person's purposes have a unique maximal member.

Understanding an action, in the chapter's selected special case, means
knowing the intention with which it was performed. The action's
subjective meaning is its function for its performer, rather than a
meaning established solely by membership in a linguistic community.
Meggle connects this explanation to Weber's understanding of subjective
meaning. His second definition compresses full understanding into
\[
 U(Y,D(X,f))\defeq K(Y,I(X,f,A^*)),
\]
where \(A^*\) represents all the performance's intended aims.
Here the source provisionally sets \(K(X,A)\defeq B(X,A)\land A\):
knowledge is replaced by true conviction for simplicity, explicitly
acknowledging that this is not an adequate general analysis of knowledge.
The compression also presupposes that the represented aims and their
combination are available; it is not an existence theorem for an arbitrary
infinite set of intentions. Its claim that understandable action is
rational is about this intention-and-reason explication, not a finding
that every action observers understand is objectively good or justified
\cite[p.~602]{Meggle2021RationalityCommunication}.

\section{Informing someone is not merely causing a belief}

Utterances are not the only communicative performances. Meggle includes
gestures, facial expressions and flying a flag; almost any performance
can be used communicatively. He concentrates on \emph{informatives}:
speaker \(S\) performs \(f\) toward hearer \(H\) with the primary
aim that \(H\) later believe proposition \(p\). The word does not
presuppose that \(p\) is true or that \(S\) is sincere.
Write \(CA(S,H,f,p)\) for the selected informative attempt. Its first
necessary condition is
\[
 CA(S,H,f,p)\to I(S,f,B'(H,p)).
\]
Note 3 gives the parallel case of a directive by replacing the target
belief \(B'(H,p)\) with a later action \(D'(H,a)\). That substitution
preserves this abstract aim/recognition structure; it does not equate
all institutional conditions for promises, commands or directives.

Causing the target belief through an unnoticed manipulation is not
enough for this account. The distinctive means is the hearer's recognition
of the attempt itself. The speaker expects the belief to arise through
that recognition, rather than merely through an unrelated cause.
The chapter combines these conditions as its adequacy requirement
\[
\begin{split}
 CA(S,H,f,p)\leftrightarrow{}&I(S,f,B'(H,p))\land\\
 &B(S,B'(H,p)\leftrightarrow K'(H,CA(S,H,f,p))).
\end{split}
\]
It calls this AC1 and explicitly refuses to treat it as a definition:
the occurrence of \(CA\) on both sides would be circular. AC1 constrains
an acceptable explication; merely writing it does not construct one
\cite[pp.~602--603]{Meggle2021RationalityCommunication}.

\begin{hbtheorem}{thm:p18-r103-reflexivity}{proved=annex:p18.r103.reflexivity}{The reflexivity consequence requires attitude principles}
Fix an actual performance \(d=D(S,f)\), target belief
\(b=B'(H,p)\), and recognition \(r=K'(H,CA(S,H,f,p))\).
Use the preceding definition of \(I\). Suppose AC1's forward direction
gives \(I(S,f,b)\) and \(B(S,b\leftrightarrow r)\).
Suppose in addition that this agent's belief set is closed under the
classical consequences used here, and that the desire bridge holds:
\[
 W(S,A)\land B(S,A\to Q)\to W(S,Q).
\]
Then \(I(S,f,r)\). Thus the source's reflexivity condition
\(CA(S,H,f,p)\to I(S,f,K'(H,CA(S,H,f,p)))\) follows under these
additional assumptions, not from propositional logic and AC1 alone.
\end{hbtheorem}

The bridge is a substantive idealized condition on attitudes. The source
later states it and cites Kant's requirement to will necessary means;
note 8 again recommends a subjunctive or strict conditional rather than
material implication. Unlimited transfer of desire through every believed
material implication can import accidental or irrelevant consequences.
The conditional derivation establishes exactly what the stated assumptions
yield, not that actual agents always satisfy them or should want every
logical consequence of an end.

\subsection{The transparency hierarchy and its boundary}

The source starts the hierarchy with
\(I_0=I(S,f,B'(H,p))\) and writes
\(I_{n+1}=I(S,f,K'(H,I_n))\) for every \(n\geq0\).
It treats a communicative intention as open to recognition at every
level, and denotes the total hierarchy by \(I^*\). Its strengthened
reflexivity condition RC-star and third definition have the forms
\[
\begin{aligned}
 CA(S,H,f,p)&\to I(S,f,K'(H,I^*)),\\
 CA(S,H,f,p)&\defeq
 I_0\land B(S,B'(H,p)\leftrightarrow K'(H,I^*)).
\end{aligned}
\]
Grice's initial model, as reconstructed by Meggle, uses \(I_0\) where
the latter display uses \(I^*\). The difference is not just notation:
the strengthened form requires openness about the intention to secure
recognition, the intention that that intention be recognized, and so on.

\begin{proofboundary}{def:p18-r103-transparency-boundary}{An infinitary explication is not a supplied finite calculus}
Each finite \(I_n\) is specified recursively. An unbounded conjunction
\(I^*\), however, is not automatically a formula of an ordinary finitary
modal language. Nor does knowing each of a list of intentions separately
automatically establish knowing their infinite conjunction in an unspecified
logic. The chapter says that its third definition entails AC1 because
knowledge of the communicative attempt amounts to knowledge of \(I^*\).
That identification needs the intended attitude semantics, relevant
introspection and closure principles, and a well-defined treatment of the
hierarchy; the printed equations alone do not supply them.

Accordingly the displays above preserve the source's proposed explication
and scope of openness, not a newly proved existence, uniqueness or
completeness result. The proved reflexivity result makes its additional
attitude assumptions explicit. Finite common-knowledge models below concern
the separate group-meaning construction and do not silently solve the
intention-hierarchy problem.
\cite[pp.~602--603]{Meggle2021RationalityCommunication}
\end{proofboundary}

The chapter's fourth definition makes understanding an informative
attempt \(U(Y,CA(S,H,f,p))\defeq K(Y,CA(S,H,f,p))\).
The interpreter \(Y\) may be the hearer or a third person. Under this
account the proposition known is that an attempt with content \(p\)
was made, not necessarily \(p\) itself. The hearer's reasons to accept
\(p\) introduce two additional normality conditions: sincerity,
\(CA(S,H,f,p)\to B(S,p)\), and correctness,
\(B(S,p)\to p\). Neither is constitutive of every individual attempt
in the selected theory.

\begin{hbtheorem}{thm:p18-r103-success-truth}{proved=annex:p18.r103.success-truth}{Recognition, uptake and truth require different conditions}
Recognition of a content-bearing attempt, even when it is known to be an
attempt to induce belief in \(p\), does not propositionally entail
\(B'(H,p)\) or \(p\). Attainment of the informative aim
\(B'(H,p)\) does not entail \(p\). Sincerity and correctness together
with \(CA(S,H,f,p)\) do entail \(p\), but these conditions alone
do not entail the hearer's uptake.
\end{hbtheorem}

An editorial contrast makes these distinctions usable. A red flag has an
established local meaning that a passage is unsafe. The observer may know
that the attendant is trying to convey precisely that message, yet decline
to accept it after independently checking that the passage is safe.
Conversely, an attendant who knows the passage is safe may use the flag
to induce a false belief; if the observer believes it, the informative aim
succeeds despite insincerity and falsity. An honestly mistaken attendant
is sincere but incorrect. These stipulated cases develop the source's flag
and sincerity/correctness examples; they are neither observed experiments
nor evidence about real safety conditions.

\section{From individual performances to shared meaning}

\begin{hbdefinition}{def:p18-r103-meaning-levels}{Four levels that must not be collapsed}
The chapter separates the \emph{performance} from its \emph{product},
and the token from the type. A particular singing, whistling or flag-flying
is a performance token (level I); its action type is level II. A song,
whistle or displayed flag considered as a product form is level III;
a particular product or utterance is level IV. One token may be an
instance of a conventionally meaningful type without its particular
speaker's purposes being identical to that type's meaning.
\cite[pp.~603--605]{Meggle2021RationalityCommunication}
\end{hbdefinition}

Methodological individualism begins with the intentions of performers
of tokens and constructs the social level from them. Methodological
collectivism starts with already meaningful action types and their rules.
This choice of explanatory starting point does not by itself settle an
ontological dispute about whether individuals or collectives exist.
Intersubjective meaning is relative to population \(P\), situation
type \(\Sigma\), action type \(F\), and content \(p\):
\(M(P,\Sigma,F,p)\). Meanings may differ between groups and between
types of situation within one group.

Meggle's fifth definition first specifies a regularity, \(M_0\): for
every speaker \(S\) in \(P\) and every token \(f\) of \(F\)
performed in a situation of type \(\Sigma\), performance implies an
informative attempt to some hearer \(H\) in \(P\) with content
\(p\). It does not require that the hearer believe \(p\), that the
speaker is sincere, or that \(p\) is true. It can hold vacuously when
there are no relevant performances; it is not by itself evidence of a
living convention. Definition 6 then sets
\[
 M(P,\Sigma,F,p)\defeq K^*(P,M_0(P,\Sigma,F,p)),
\]
where \(K^*\) is the chapter's common-knowledge notation: everyone
knows the regularity, everyone knows that everyone knows it, and so on.
This construction defines type-level meaning \emph{in terms of} token-level
attempts; it does not identify conventional meaning with a particular
speaker's meaning. The source criticizes that identification, including
the invitation to it in Bennett's meaning-nominalist terminology.

\begin{hbtheorem}{thm:p18-r103-common-meaning}{proved=annex:p18.r103.common-meaning}{The common-knowledge adequacy condition under specified semantics}
Use the existing finite-agent epistemic semantics of
\cref{def:p08-rat26-epistemic-update,entry:p08-common-knowledge}.
For a nonempty finite group \(P\), identify the source's \(K^*(P,q)\)
with the group modality \(C_Pq\), interpreted over the
reflexive-transitive closure of the agents' accessibility relations.
For a well-defined regularity proposition \(q=M_0(P,\Sigma,F,p)\),
put \(M=C_Pq\). Then
\[
 M\to C_PM,\qquad M\to q.
\]
The first is the source's AC2, \(M\to K^*(P,M)\); the second records
the factivity needed to infer the regularity itself. These are conditional
semantic facts, not a proof that a given community realizes the model.
This interpretation supplies explicit accessibility semantics; it is not
entailed merely by the earlier abbreviation \(K(X,A)=B(X,A)\land A\).
\end{hbtheorem}

Bare shared regularity still leaves a rationality gap. A speaker expects
both recognition of the attempt and acceptance of its content. Commonly
known meaning helps explain expected recognition, but the proposed
definition of \(M\) contains no general reason for repeated acceptance.
Meggle appeals to Lewis's convention account: shared interests together
with mutually known strategies can stabilize regularities. His ideal
\emph{perfect communication situation} requires three distinct ingredients:

\begin{enumerate}
\item A common interest that the hearer's later belief tracks truth,
      \(B'(H,p)\leftrightarrow p\).
\item Common knowledge that the speaker knows whether \(p\), while
      the hearer initially knows neither \(p\) nor its negation.
\item Common knowledge of a joint strategy: the speaker performs
      \(f\) exactly when \(p\), and the hearer believes \(p\)
      exactly when that performance occurs.
\end{enumerate}

\begin{hbtheorem}{thm:p18-r103-perfect-signal}{proved=annex:p18.r103.perfect-signal}{Truth transmission in the ideal joint strategy}
Let \(d=D(S,f)\) and \(b=B'(H,p)\). If the actual strategy conditions
\(d\leftrightarrow p\) and \(b\leftrightarrow d\) hold, then
\(b\leftrightarrow p\). If moreover the speaker knows either
\(p\) or \(\neg p\), and knowledge is factive and includes belief,
then an actual performance \(d\) entails \(B(S,p)\) and \(p\).
This establishes sincerity and correctness for the observed signal
under the ideal assumptions. Common knowledge of the strategies supplies
the actual conditions only when its semantics is factive.
\end{hbtheorem}

Common interest explains why parties might adopt the arrangement; it is
not needed for the displayed truth-functional implication once strategies
are assumed. Likewise common knowledge supports expectation but does not
prove that incentives, monitoring or repeated learning actually maintain
the arrangement. In the absence-of-signal case the model ensures absence
of positive belief \(B'(H,p)\), not necessarily positive belief in
\(\neg p\): that further attitude is not one of its strategy clauses.

Note 4 explicitly says that the universal quantifiers describe an ideal
type and proposes weakening them to high percentages, following Lewis's
later suggestion. The chapter supplies no percentage, denominator,
population sample, interaction task, design, estimate, uncertainty interval,
learning dynamic or comparative test. Its statement that sufficiently
frequent confirmation supports stable conventions is a qualitative
explanatory proposal here, not a demonstrated universal collapse threshold.
No transfer to adversarial speakers, different communities, online media
or unequal-power interactions is licensed by this ideal example
\cite[pp.~603--604, 608; p.~609, n.~4]{Meggle2021RationalityCommunication}.

\section{Semantics, pragmatics, and constitutive rules}

For a simple unstructured sign, such as a light or flag, the source
identifies its product-form meaning with the conventional meaning of
producing it. For structured expressions this shortcut leaves the
compositional problem open: how meanings of parts determine meanings
of wholes. Meggle describes an action-theoretic semantics that would
interpret a syntactically specified language enriched with indexicals,
intensional and performative devices by communicative conventions.
He refers to Lewis, von Kutschera and his own work for selected formal
developments; extending such a construction to natural languages is an
additional task, not a theorem contained in this chapter. His assessment
that no generally accepted compositional answer is available is a
source-attributed judgment, not a proof that such a semantics is impossible.

At level IV, ordinary pragmatics concerns what concrete utterances do
in their actual situations. A token's belonging to an ideal situation
type does not exhaust its circumstances. Literal and nonliteral uses,
speaker goals, institutions, audience knowledge and cultural practices
can alter the explanation of expected understanding and uptake. The
organizing question is which features justify those expectations in this
particular exchange. The chapter places Gricean implicatures, relevance
theory and conversational maxims in this wider field but expressly leaves
their detailed procedures beyond its scope
\cite[pp.~604--605]{Meggle2021RationalityCommunication}.

Speech-act theory, as Meggle presents it, starts with intersubjectively
meaningful act types: commanding, asking, questioning, thanking,
threatening, warning, recommending, dissuading, declaring, excusing,
promising and others. These illocutionary roles
form an open list, not a complete finite classification supplied here.
Meggle criticizes the absence of an explicit general definition beyond
such lists and regards conceptual explication as an outstanding need;
that is his assessment of the tradition, not a census undertaken here.
Because conventional linguistic signs are also presupposed, the starting
point combines levels II and III, rather than treating every nonlinguistic
communicative token from the outset. Note 5 recalls Austin's locutionary
meaning in the Fregean sense. Meggle discusses Alston, von Savigny and
Brandom as differing attempts to connect structured sentence meaning
with pragmatics, not as interchangeable versions of one reduction.

The central irreducibility thesis says that illocutionary acts or forces
cannot be defined by perlocutionary acts or effects. Perlocutionary acts
bring something about through doing or saying something; intentionally
doing so fits the earlier instrumental account. Meggle's objection
distinguishes two kinds of illocution. Institutional acts, illustrated
by marriage, baptism and the source's auction example, can count as the relevant acts
because the procedure's constitutive conditions hold, without one special
instrumental aim being conceptually required. A person may seek wealth,
social acceptance or redemption, but none of these purposes alone defines
the institutional act. Meggle treats irreducibility as valid for this
kind of illocution on that conceptual ground.

Communicative informatives present a different case. The individualistic
explication analyzes them using a special open, recognition-mediated
instrumental intention. Meggle therefore rejects a blanket extension of
institutional irreducibility to those communicative acts. This is his
reductionist proposal and argument, conditional on accepting its
explication; it is not an empirical experiment establishing that every
illocution has that analysis. A collectivist can instead refuse the
reduction and characterize meaning by the constitutive rules for correct
performance of a type in a population and situation. Understanding then
means knowing those rules, rather than knowing the particular performer's
intention. Note 6 distinguishes the role of language-game rules from
chess rules and cites further discussion
\cite[pp.~605--606; p.~609, nn.~5--6]{Meggle2021RationalityCommunication}.

Correctly following chess rules does not explain why someone wants to
play now or why this legal move is preferred to another. The source uses
that example to show why constitutive correctness leaves questions of
action rationality alive. Correspondingly a valid institutional declaration
does not, merely by its procedural standing, establish the truth of an
unrelated factual allegation. Truth, institutional standing, recognition,
and perlocutionary effect must be assessed with their respective warrants.

\section{Habermas's different collective question}

In Meggle's reconstruction Habermasian communicative action belongs to
a collective, not merely to a single member making an informative attempt.
Individual contributions participate in a common practice with a shared
attitude. The source separates three parts of that attitude:
\emph{cognitive}, common recognition that a factual, expressive or normative
issue is open; \emph{voluntative}, a common interest in open-minded pursuit
of mutual agreement; and \emph{praxiological}, giving and considering
reasons through theoretical or practical discourse.

The speech-act background has objective, social and subjective dimensions,
associated with representation, appeal and expression in the discussed
B\"uhler tradition and with culture, society and personality. Corresponding
validity claims are to be redeemed discursively. The source links this
program to Apel's transcendental pragmatics and the Peircean idea of
long-run agreement by an indefinitely extended community. These are
historical and philosophical programmatic connections, not a proof that
any actual consensus is true, or that one procedure resolves every factual
and normative dispute. The chapter expressly says it cannot reconstruct
Habermas's much larger multiperspective research program in this space
\cite[pp.~606--607; p.~609, n.~7]{Meggle2021RationalityCommunication}.

Meggle's comparison retains four material differences. First, the
irreducibility thesis is a substantive foundation to be examined, not a
neutral bridge already accepted by both approaches. Second, Habermas's
openness concerns a jointly unsettled issue \(p\) or \(\neg p\).
The individual informative's openness concerns the fact that the speaker
is making the attempt; it need not include open-minded reconsideration
of \(p\). Third, Habermas's opposed category of strategic action lacks
the relevant openness by definition, while strategic pursuit of an end
can remain communicative in the individualistic account. An openly
recognizable but insincere assertion is not eliminated by that account's
definition. Fourth, collective agency is \emph{sui generis} at the level
of its declared subject, but calling it collective does not prove it is
incapable of explanation by individual attitudes and practices.

The reduction question is thus left open for the broader collective
program. Meggle concludes that an overall analytic reconstruction and
connections to developed theories of collective intentionality would be
valuable; his account of the state of that literature is his 2021
assessment, not a measured exhaustive census or a finding about work
published subsequently. He identifies Steinhoff's more critical and
Heath's more constructive work as promising beginnings, without treating
either as the completed overall reconstruction
\cite[pp.~607--609]{Meggle2021RationalityCommunication}.

\section{Three rationality questions and their interfaces}

The final source section distinguishes action rationality (the action's
fit to the agent's beliefs and preferences), rationality of its reasons
(whether those attitudes and their interrelations satisfy the chosen
standards), and personal rationality (whether the person actually acts
on the reasons available in the situation). Having adequate reasons
does not itself entail acting on them. Belief consistency, deductive
closure and desire consistency are possible idealized constraints,
not observations of unlimited human inference. Relations across time
also require a theory of revision; static consistency does not provide it.
The desire bridge used in the reflexivity proof is one explicitly
identified cross-attitude condition.

If speech-act theory accepts the proposed individualistic reconstruction
of informatives, its rationality assumptions for that class coincide
with those of the reconstruction. If it retains constitutive rules and
the irreducibility thesis instead, understanding has a different object.
The chess example shows why that version still needs an account of the
agent's reasons for performing one permitted act rather than another;
constitutive correctness is not a complete explanation of practical choice.

In the individualistic approach, communication adds special intended
openness and expectations about the hearer's interpretation, rather than
a wholly new standard independent of action rationality. The soundness
of a speaker's reasons concerns objective success and warrant, while the
individual action can remain subjectively explicable under false beliefs.
At the conventional level, mutual recognition and reliably confirmed
uptake supply further reasons for expectations. Game-theoretic modeling
here concerns a collective \emph{distributively}: the decision makers
remain individual members. That is not the same as Habermas's collective
being the subject of an action. The source's reference to common knowledge
of rationality in note 9 must also not be changed into common \emph{belief}
with no truth condition
\cite[pp.~608--609, nn.~8--9]{Meggle2021RationalityCommunication}.

That subjective and intersubjective support also permits communication
between the lines. Meggle treats Grice's cooperative principle as basic
within the implicature theory but as derivable from a cooperative
rationality perspective. This is an explanatory relationship asserted
in the selected account, not a fully developed implicature calculus here
or an assumption that every actual speaker is cooperative.

These distinctions have specific handbook interfaces. Hahn and Collins's
procedural, epistemic and conventional standards
(\cref{def:p18-r55-norms}) separate orderly participation from evidential
support and empirical conformity to a norm. Meggle's \(M\) is instead
a proposed common-knowledge construction of meaning; his use of
\emph{conventional} is not their empirical validation query. Woods's
consequence-having, spotting and drawing
(\cref{def:p18-r56-three-consequence-roles}) likewise separates a static
inferential relation from an agent's recognition and action. Knowing
what a person tried to communicate supplies no automatic warrant for
the content's truth.

Walton's actual dialogue-purpose distinctions
(\cref{def:p18-dialogue-purpose}) retain inquiry, negotiation, persuasion,
information seeking, deliberation, discovery and eristic episodes. A
negotiated settlement is not by that achievement a factual discovery;
identifying an informative attempt does not identify the entire dialogue
as a Habermasian open-minded inquiry. The P11 attitude-based distinction
(\cref{def:p11-rat22-attitude-types}) prevents confusion of a belief
about an action with the practical attitude of deciding to perform it.
The finite epistemic-game account
(\cref{def:p17-egt-types,def:p17-egt-common-belief}) explicitly permits
mistaken beliefs. Its common-belief condition cannot substitute for
the factive common-knowledge premise used in the ideal signaling proof.
Conversely the P8 truthful public-update model
(\cref{def:p08-rat26-epistemic-update}) builds truth and public observation
into a protocol; an arbitrary sincere or insincere informative is not
automatically that operation.

\begin{hbmethod}{def:p18-r103-assessment}{Assess the actual communicative commitment and its warrant}
Identify the performance token, speaker, addressed hearer, content, time
and purpose. Distinguish a communicative attempt from an unnoticed cause
of behavior. State which kind of understanding is at issue: recognition
of the performer's intention, mastery of a meaningful type's constitutive
rules, or participation in a collective inquiry toward agreement.

Separate intended uptake, actual uptake and truth. For information offered
by a speaker, examine sincerity, competence, access to the relevant facts,
and independent support; understanding the assertion is not evidence
that all these conditions hold. For an institutional act, identify the
actual constitutive conditions instead of replacing them with a presumed
private purpose. For a consensus, state the issue, openness and reasons
considered; do not infer truth solely from agreement or procedural success.

When using conventions, fix the population, situation type and action
type, then distinguish observed regularity, mutual knowledge and incentive
support. A flag used once with recognizable intent is not evidence of a
stable convention across populations. An ideal joint-strategy example
can demonstrate a conditional implication but supplies no measured
frequency or warranted policy for a noisy or adversarial setting.

Finally state which verdict was actually obtained: successful recognition,
successful persuasion, reliable information, correct institutional
performance, an accepted practical choice, or justified agreement.
Failure of one verdict need not be failure of all the others. This is an
editorial assessment procedure grounded in the chapter's distinctions,
not a complete source-defined decision algorithm.
\end{hbmethod}
